% DECISION: the MLSys'27 call says to upload the appendix separately from the main paper, at the same deadline, with
% no length limit. This file is therefore its own document, not a section of main.tex, and it builds independently.
% DECISION: nothing here is load-bearing. Every item has a sentence in the body that stands on its own, and this
% document carries only the detail a sceptical reviewer would ask for. A claim that lives only here is a bug.
\documentclass{article}

\usepackage{microtype}
\usepackage{graphicx}
\usepackage{booktabs}
\usepackage{array}
\usepackage{amsmath}
\usepackage{amssymb}
\usepackage{textcomp}
\usepackage{xcolor}
\usepackage{threeparttable}
\usepackage[hyphens]{url}
\usepackage{xr}\externaldocument{main}
\usepackage[hidelinks]{hyperref}

\newcommand{\theHalgorithm}{\arabic{algorithm}}
\usepackage{mlsys2025}

\newcommand{\E}[1]{\textcolor{red}{$\langle$#1$\rangle$}}

\newcommand{\us}{\,\textmu{}s}
\newcommand{\ms}{\,ms}

\let\oldtexttt\texttt
\renewcommand{\texttt}[1]{{\small\oldtexttt{#1}}}

% The body refers to these sections as Appendix A, B, ... since the appendix is uploaded as its own document.
\renewcommand{\thesection}{\Alph{section}}

\mlsystitlerunning{Appendix: OPCORD}

\begin{document}

\twocolumn[
\mlsystitle{Appendix: OPCORD: Recovering Operator Concurrency from Command Order on AMD AI Engines}

\begin{mlsysauthorlist}
\mlsysauthor{Anonymous Authors}{anon}
\end{mlsysauthorlist}
\mlsysaffiliation{anon}{Anonymous Institution}
\mlsyscorrespondingauthor{Anonymous Authors}{anonymous@anonymous.org}

\mlsyskeywords{AMD AI Engine, XDNA 2, LLM inference, command processor, scheduling}

\vskip 0.3in
]

\printAffiliationsAndNotice{}

\section{Attributing an Issue to a Column Group}
\label{app:attribution}

% VALUE: the one disclosure in the paper that a reviewer should be able to audit in full: the attribution rule was
% settled after the capture, the earlier rule labels the same capture differently, and the evidence that the later
% rule is the physically right one comes from waits and argument slots rather than from the rule itself
Section~3.4 of the paper attributes each issue in a recorded command stream to the column group whose compute tiles its
data reaches, rather than to the column the instruction addresses. This appendix states how that rule was chosen,
because it was settled after the capture it is applied to and the earlier rule gives a different label.

\textbf{Why the addressed column is not the group.} In the production engine's fused prefill, each attention command
also pushes one shim queue on a column belonging to the other group: one image pushes on column~6 and another on
column~3. Attributing those pushes by the column they address places them in the GEMM group, inside commands that are
otherwise attention, which yields 10{,}822 group transitions of which 10{,}752 are internal to a command.

\textbf{Why those pushes belong to attention.} The evidence is independent of the attribution rule. The two images'
synchronizations are on the attention column only: 128 waits all on column~7 for the first, 256 all on column~0 for
the second, and no GEMM image touches column~0 or~7 at all. The borrowed channel reads a different argument slot than
the attention column's own while carrying that slot's exact per-block pattern, which is a second input streamed in
lockstep rather than GEMM work. The route itself is in the configuration, not in the control code.

\textbf{The two labels.} Under the rule registered first, by addressed tile column, the capture is mixed, with 99.4\%
of its transitions internal to a command. Under the rule the paper uses, where each issue takes the group whose compute
tiles its data reaches by the configuration's routes, the capture has 70 transitions, 35 into attention and 35 out of it,
and every one sits at a boundary between two commands, which is the serial label. We record plainly that the second
rule was proposed after the capture and moves the outcome toward our claim: under the rule registered first, the
compute-weighted share of the transitions of the overlapping kind is 0.776 (its bytes-only share 0.912), which is the
majority-mixed label, so the serial prediction we registered was a miss under the registered instrument and a hit
only under the rule chosen afterwards.

\textbf{What a population check adds, and what it does not.} Over all 218 recorded streams, the columns that carry any
instruction fall into four sets: 200 streams address only the GEMM columns, 8 address columns $\{0,3\}$ and 8 address
$\{6,7\}$, 1 addresses all eight, and 1 carries no instruction. In the capture itself, 276 of 311 commands address only
the GEMM columns, 28 are one image on $\{6,7\}$ and 7 the other on $\{0,3\}$, and the stream that addresses all eight
is not dispatched. In every two-column stream the bulk of the instructions sit on the attention column, by 1{,}420
against 384 in the largest case. This counts addressed columns and not the programs those instructions serve, so it
corroborates the attribution rather than establishing it; and because a rule that labels each command by the
columns it addresses also gives 70, the fact that the column set changes exactly 70 times re-derives that count
rather than measuring it a second time.

\section{Other Round-Trip Populations}
\label{app:roundtrip}

% VALUE: the lead contribution's number will be audited, so the populations beside it are stated rather than left for
% a reviewer to find
The 625 streams quoted in the introduction are the streams each toolchain emits for a design. Three further
populations round-trip byte-identically and are not counted there, to keep that number tied to one kind of object.
The open compiler also dispatches streams that differ from its on-disk files, either stripped to a column group or
interpolated from a template: 539 of those re-emit byte for byte, and none of them is among the 375, so that toolchain
has 914 distinct streams verified across two gates. Our own composer's output re-emits on 2{,}684 streams, which we
exclude from a toolchain count because it is our output rather than a stack's. The production engine's streams at
other prompt lengths, 37 and 53 in two further gates, overlap the 218 already counted. Each gate carries two negative
controls, a modified stream and a truncated one, which fail as designed.

\section{The Production-Engine Comparison}
\label{app:flm}

% VALUE: the inequalities behind the reference point, and an honest statement that we can price only part of the gap
The comparison with the production engine in Section~5.2 is a reference point, not a like-for-like speedup. The two
stacks differ in their kernels, their dequantization path, their tokenization, and the chunk size their prefill uses,
and we hold the weights and the prompt fixed but not those.

\textbf{What we can price.} In a capture of the engine's dense path at two prompt lengths, the command counts are
fixed: 146 commands, 64 configuration switches, 12 module and kernel constructions and 226 runs per request at both
lengths. A separate dequantization pass is fixed in count and in time, 16 run lists of four commands taking
78--84\ms{} at both lengths, and one further pass takes 18--20\ms. That is about 100\ms{} of fixed device time, which
is 40\% of the 248\ms{} difference from the shipped program at 2{,}048 tokens and 49\% of 208\ms{} at 4{,}096. Every command in the capture is
serial, and the unit repeated per layer is a dequantization, three matrix multiplications, an attention command and
four more matrix multiplications.

\textbf{Where the difference goes.} Recording both stacks the same way, by the time each command is outstanding on
the device and the time the host spends with nothing outstanding, accounts for the whole of the difference at a
4{,}096-token prompt, where the difference from the shipped program with attention alternated was 194\ms{} on the days of the two recordings. Host time accounts for more than all of those 194\ms{}: the production engine spends 318\ms{} with no command
on the device against our 96\ms. Its fixed passes account for a further 102\ms. Its matrix work runs 65\ms{} longer
than ours. Against those, its attention runs 196\ms{} \emph{faster} than ours, which is the one place the hand-tuned
engine is ahead of the design we rewrite, and which is why the difference is not larger.

Two limits on that accounting. The comparison is by command kind rather than by identical work: our matrix figure
includes the LM head and the dequantization folded into each multiplication, which the engine pays
elsewhere. And the recording locates the host time without seeing into it, since the interposer observes the gaps
between calls rather than what runs inside them, so the dominant term is host-side time we can place but not
decompose. At a 2{,}048-token prompt we have no equivalent record of the open stack and attribute nothing.
% SOURCE: vault lab-flm-streams (15), no device. FastFlowLM from the capture's serial issue-to-done per kind; the open
% stack from the ML Timeline hold-2 records, 1B shipped with attention alternated at 4,096 tokens, TTFT 1,607.6 ms,
% per-command first issue to last return unioned per column group, median of 5 reps: GEMM group 895.0 ms, attention
% group 617.1 ms, overlap 0.0, host time with nothing outstanding 95.6 ms. FastFlowLM: GEMM 959.9, attention 420.8,
% separate dequantization and layer passes 102.5, host idle 318.5. Differences +64.9, -196.4, +102.5, +222.9, summing
% to +194.0 against a measured window difference of 194.1 on these two days. The two sides come from different holds
% on 2026-10-21.
\section{Registered Predictions: Flagged Arms}
\label{app:flags}

% VALUE: the arms whose spread or bracket exceeded its threshold, listed so that a reader can see that none is a
% composer arm and that all were scored as they stand
Four baseline arms varied by more than 5\% across their repetitions: the shipped program at 12{,}288 + 2{,}048 tokens on
one split of Llama-3.2-1B (+6.3\%), the shipped program at 2{,}048 + 4{,}096 on another split of the same model
(+6.7\%), and at 2{,}048 + 4{,}096 on Qwen3-4B both the shipped program (+8.4\%) and the shipped program with its
attention columns alternated (+5.8\%). One bandwidth bracket, on a Gemma3-4B split that the composer does not choose,
ran 7.3\% above the day that split's probes were taken; it touches only comparisons at that split, and the split the
composer chooses for Gemma3-4B is the other one throughout. No composer arm was flagged. Every flagged arm is scored
as it stands, and none was re-run after its outcome was seen.

\section{The Synthetic Designs}
\label{app:synthetic}

% VALUE: the three disclosures that belong with the synthetic measurements, including the one hold whose bandwidth
% bracket is missing and an observation the registration did not anticipate
Three disclosures belong with Section~5.4's synthetic designs.

First, in the all-or-nothing design at the smallest level of core work, the repetitions of the depth-one baseline span
193\% of their median, because that variant was the first opened in the hold; its median is the statistic the
registration names, and the other spreads in that design are between 0.4\% and 4.9\% except three, of 34\% at the same level and 15\% and 14\% at the middle one.

Second, the bandwidth cap probe did not run for that hold, because a relative output path was passed to a probe that
changes directory first, so no cap reading brackets those measurements. It is not one of the output checks, all of
which passed. The other two holds read 13.4--13.5~GB/s at one column and 54.7--55.1~GB/s at eight columns at both ends.

Third, an effect we did not anticipate and can now account for: in the order-by-loop
design at buffer depth four, both orders run at 11.65\ms, 3.2 times the depth-two time of the faster order, where a first
reading of the period had expected the two depths to agree. We did not alter the design after measuring it. No cycle of the depth-two graph accounts for it, since adding buffer slots can only add tokens to a cycle
and cannot raise any cycle's mean, so the firing times themselves differ at depth four. Comparing the final
configuration state of the four built designs localises that: every shim-tile and memory-tile register is identical
between the depths, with no change to a transfer, a route or a lock outside the cores' own tiles. Reading the lowered
code names it. The buffer selection lowers to a two-way choice between constants at depth two and to a load from a
three-entry table of addresses at depth four, and the inner loop over 1{,}024 elements is unrolled at both depths but
software-pipelined only at depth two: 12 instruction bundles per four elements against 38, with each load at depth four
waiting out its latency. Those counts predict 6.29 and 19.92 million cycles, which match the measured 3.610\ms{} and
11.653\ms{} at 1.74 and 1.71\,GHz against an 1.8\,GHz clock, so the bundle counts account for the slowdown to within
2\%. Why pipelining stops is not established; a pointer loaded from a table cannot obviously be proven not to alias the
input, which would block hoisting loads above stores, but we did not check that, and we have not excluded the memory
banks the depth-four buffers share. What the certificate compares is unaffected, since both orders at depth four run
the same loop; what changes with depth is the time of that loop, and the mechanism for it is named above.

\section{Reference Checks and the Gemma3-4B Gate Result}
\label{app:gemma-gate}

% VALUE: the three instruments differ in kind, which matters more than how many tokens each covers; then the 1B
% result that clears the headline model's prefill path, then the Gemma result with its scope
Three kinds of instrument are available to us, and they are not interchangeable. The expected continuations shipped
with the Qwen3-4B and Gemma3-4B modules are constants their authors recorded from their own device runs, one of them
re-recorded after the authors fixed a bug that had zeroed a feed-forward contribution; agreement with such a constant
is a regression check against the authors' output over ten tokens, not a statement of numerical fidelity. The
remaining instrument is a floating-point forward on the same quantized weights, which two of the stacks provide and
which the designs are themselves gated against; comparing to it detects drift of the device's arithmetic from that
forward, and it is the instrument used for the two checks below. It is not an independent implementation of the
model, and we do not claim it as one.

\textbf{Llama-3.2-1B.} Against that forward, the shipped prefill's keys differ by 0.4--3.0\% and its values by
1.1--7.5\% across the sixteen layers, on a six-token prompt and on a 256-token one, and the two split builds give the
same errors to every digit we report. The error grows over the first eight layers and then holds, with no step of the
kind the next paragraph describes. The first token matches the forward's on both prompts and both builds. Seeding the
turn's decoder from the forward's own keys and values, from one split's, and from the other's produces the same
seventeen tokens in all three cases, so these differences do not change what decode emits. All three agree with the
forward's own continuation for ten tokens and then differ, at a step where the forward's margin between its top two
tokens is 0.16, its second-narrowest in that sequence; because the difference appears even when the decoder is seeded
from the forward's own keys and values, it lies between the device's decode and that forward rather than in the
prefill we rewrite.

\textbf{Gemma3-4B.}

% VALUE: the per-layer detail behind the body's statement that one configuration fails a reference check, with the
% scope stated so that it cannot be read as either our build being broken or the stack being inaccurate
Section~5.1 reports that we measure no turn on Gemma3-4B because the configuration failed a reference check before any
run was timed. The check crossed two sources of keys and values, the module's own reference forward and the path the
turn harness uses, with two decode paths, on a six-token prompt with nine greedy dispatches.

Both reference-sourced arms reproduce the reference continuation exactly, and for either source the two decode paths
give logits that agree bit for bit, so both decode paths and the seeding step are clean. Both arms sourced from the
split path diverge at the fifth decoded token, with the same continuation the turn run produced. The keys
differ from the reference by 0.6--1.8\% in layers 0--9 and 5.0--9.1\% in layers 10--33, and the values by 0.8--4.9\% in
layers 0--31 and by 6.8\% and 13.0\% in the last two. The divergence lands on a near tie: with the reference's keys and
values the correct token leads by 0.875, and with the split path's the other leads by 0.25.

Two further measurements bound what this result is. The two floating-point forwards available for this model, the
one used above and the stack's other, agree with each other to 0.3--1.1\% on keys and 0.2--1.7\% on values, and give
the same first token with the same five leading candidates in the same order, so the difference reported here is five
to nine times the disagreement between the instruments and is not an artifact of which one served as the oracle. It
is also not of a different order from this stack's usual quantization error: on Llama-3.2-1B the shipped prefill
differs from its own forward by up to 3.0\% on keys and 7.5\% on values without changing a single decoded token. What
distinguishes this case is that the larger difference fell on a near tie.

The keys and values that fail came from the shipped prefill program, with no rewrite of ours applied, and were
stacked by the compiler's own routine. Two comparisons settle where the difference comes from. The build our harness
uses is the module's own program: rebuilt by the module's own builder, in the same checkout and without opening the
device, 22 of its 38 files are byte-identical and the remaining 16 are identical in size with their differences
confined to header identifiers and a timestamp, which do not reproduce between builds. The same holds for the two
models that carry our decode and turn results: rebuilding the baseline prefill behind every reported run of
Llama-3.2-1B, at both window sizes, and of Qwen3-4B gives 18 of 30 files byte-identical and the 12 remaining of equal
length with byte-identical images, differing only in those same identifiers. So the program we compare against is the
module's own on every model we report. This covers the baseline split alone, since the other two splits are
constructions of ours and have no counterpart to rebuild. And the keys and values do not
depend on our host path: the in-tree build loaded by the module's own path, the split build loaded by that path, and
the split build loaded through our turn harness with its specification patch, its weight arenas and its template
produce keys and values bit-identical to one another, with a maximum absolute difference of zero, each reproducing
the error above to the digit. The first attempt at this comparison did not run: its own floating-point reference step raised the machine's
load average above the threshold the protocol requires before the device is opened, and it exited at that check
before any device was opened; the comparison was repeated after a wait for the load to settle.
The difference is therefore the shipped program's own. We claim no cause for its step at
the tenth layer, and we report this model's prefill results as measured, taken on that program and producing its own
keys and values. The comparative claims elsewhere in the paper are untouched, because they compare a rewrite with the
shipped program on one build.

\section{Choosing the Split Before Reordering}
\label{app:split-first}

% DECISION: relocated from the introduction's list of reasons the loss survived (decisions log D19g, 2026-10-22): under the stronger baseline the existing program's faster split agrees with the composed one in 8 of 8 W1 cells and 11 of 12 Table 4 cells, so it is a narrow finding rather than a reason; kept whole, not deleted.
Tuning the array's partition before reordering issue rarely leads somewhere other than the partition the reordered program wants. Without the hand alternation, the faster split for the shipped program is not the one the reordered program wants at three of the configurations we measure; interleaving the shipped program's own attention columns by hand removes those three, but on Llama-3.1-8B, for a 4{,}096-token append to a 4{,}096-token context, the interleaved program still runs faster with four attention columns where the reordered one wants two.
% SOURCE: the un-alternated shipped program prefers S4 while every reordered version prefers S2 at 1B scratch 2,048 / 4,096 and Qwen3-4B scratch 4,096 (5-rep w1 files, vault paper-cosched-claims-v2, lab-split-sweep (18)); Llama-3.1-8B 4k + 4k: shipped S4 10,909 ms, shipped + Ar S2 11,440.1 ms, shipped + Ar S4 = shipped S4 (the alternation is a no-op on S4), composer S2 7,395.4 ms (lab-split-sweep (14), (20)). Corrected in the 2026-10-22 foundation audit (B14): the earlier clause said the two agree once alternated.


\section{Registered Predictions and Their Scoring}
\label{subsec:tests}

% VALUE: every prediction is tested as registered before its run: the bound against every measured removal, the event model's times, the joint prediction, and the composite test with its expected outcome stated beforehand; counts reported in full
Every quantity this section computes is predicted before the run that measures it, in a registration committed to the repository, and scored afterwards by a program fixed at the same time.
Four predictions test the event model: which split is faster at each configuration, the sign of every pairwise comparison between schedules, the absolute error of each predicted time against a tolerance stated in advance and against a linear interpolation between measured configurations, and the end-to-end gain.
Two test the bound: that no measured gain exceeds the bound computed for it, and that the composed schedule stays within a stated distance of the bound.
One tests the renaming of the decode loop, as a band on the per-token gain at two context lengths rather than a point.
One tests the set property directly, on designs written to hold two name dependences at a known position: that a removal off the binding cycle is worth nothing at every level of core work, with that outcome, and not a gain, stated as the prediction.
Section~\ref{sec:eval} reports each outcome with its counts, including the three that failed and the direction of each failure.
% DECISION: the outcomes and the verbs they allow follow outline v9 section 3's rows; a test not completed is reported as its least favourable outcome.
% SOURCE: the inventory is the set actually scored: the event model's four (Section 5.2, E1-E4, with E4 failed), the bound's two (Section 5.7, the 33-run check held and the 1.30x range failed on 5 of 12), the decode band (Section 5.5, both cells in band on Llama-3.1-8B), and the synthetic certificates (Section 5.3, held at two of three work levels). Nothing else in the paper is registered, and this list is re-checked against the scores at the results freeze.


\section{Discussion}
\label{app:discussion}

% MERGED from the former discussion section
% VALUE: transferable insights only, each a possibility under its condition and tied to a measured result: an in-order issuer rations concurrency so its order is the schedule and packing yields little of it; the faster partition may change with the regime, so evaluations of prompts from scratch only may choose the wrong one; probes may rank configurations reliably before they run; once ordered, the bound may move from the issuer to the groups; a false dependence inside an operator is out of reach of reordering but not of renaming; on agent turns, gains measured from scratch may misstate end-to-end gains
Although we evaluate one NPU, our findings rest on properties that other spatial accelerators may share.
Where an in-order issuer mediates the concurrency of an array, its instruction order may act as the schedule, so packing work into fewer dispatches may yield little concurrency, and a second queue or a completion signal per group of columns could remove the rationing in hardware (Sections~\ref{subsec:semantics} and~\ref{subsec:ablation}).
Where operators are partitioned spatially, the faster partition may change between computing a prompt from scratch and appending to a cached context, so an evaluation of prompts computed from scratch alone could choose the partition that is slower for agents (Section~\ref{subsec:ablation}).
Where issue returns at once and waits block, probes of each operator alone may rank configurations reliably even where they misprice them, so a configuration can be chosen before it runs (Section~\ref{subsec:accuracy}).
Once the issue order is right, the binding term may move from the issuer to the groups' own busy time, so further gains could come from kernels and from the balance of work between groups rather than from scheduling (Section~\ref{subsec:bound-eval}).
Where an operator reuses one buffer across the iterations of a loop, reordering issue across operators cannot reach the resulting wait, whereas renaming the buffer inside the operator could (Section~\ref{subsec:decode-eval}).
Finally, where agents append to long contexts, prefill becomes bound by attention and decode by the KV cache, so gains measured on prompts computed from scratch may misstate what an agent's turn gains (Section~\ref{subsec:turns}).
% DECISION: each sentence is re-checked against its section's result at the results freeze; a sentence whose result does not hold is dropped (vault paper-cosched-plan, section 6).

% The former limits paragraph is distributed to the points it qualifies: the passes' scope to the passes, the
% host state to the turn result, the reference-point caveat, the gates, the composition cost and the two failed
% predictions to the evaluation, and the analysis's refusals to the analysis. Only the one-device limit, which
% qualifies the transfer argument above, stays here.
We measure one generation of one vendor's array, so the transfer above rests on the properties we name rather than on measurements elsewhere.

GPUs once lost concurrency in the same way, where streams that shared one hardware work queue could wait on each other, and NVIDIA's Kepler GPUs added hardware queues that remove these false dependences~\citep{nvidia-kepler-tuning}.
On XDNA~2, one in-order CP per hardware context leaves this job to the program, so the compiler owns concurrency across column groups, much as Rammer's compile-time schedules do on GPUs and IPUs~\citep{ma2020rammer}.
This places XDNA~2 between GPUs, whose hardware decides overlap at run time, and statically timed machines such as the TSP, whose compiler fixes when every instruction executes~\citep{abts2020tsp}, because its program fixes the order of each engine's work but not its duration.
A device that reads busy for 92--97\% of TTFT while only one column group works at a time shows why this ownership matters, since utilization measured per device hides the idle columns (Section~\ref{subsec:losses}).
Other arrays whose engines run their programs in order, such as other AIE generations, may share this property, so the same analysis could apply where their streams can be read.
% SOURCE: Kepler Tuning Guide section 1.4.2.2 (read 2026-10-21); device busy = summed isolated operator times less the dispatch floor, 92/95/97% of TTFT at 512/2k/4k (amd-npu e3acc8a, Figure busy), not a hardware counter.
% DECISION: "the compiler owns concurrency on static hardware" is cited (Rammer, TSP) and kept as discussion, never a contribution (outline v9 section 2, "Not contributions").

% VALUE: answers the obvious question of why one of the three passes runs on the host while the other two are
% compiler passes. The answer is a condition rather than a boundary: the two inside the compiler are the two whose
% output does not depend on the prompt, and a deployment whose shapes are known could emit the composed stream
% ahead of the run as well.
Where this work would move into the compiler is decided by when the shapes are known.
Two of our three passes already run there, and they are exactly the two whose output does not depend on the prompt: the stream pass reorders the segments within one operator's stream, and the loop pass renames the buffers of decode attention, and one build of each serves every context length.
The composer is the piece that runs at dispatch, and it is the piece whose plan is a function of the cached context and the appended tokens.
Nothing in that plan requires a host.
Planning a schedule and composing its commands are compile-time work in kind (Section~\ref{subsec:setup}), and a compiler given the shapes could emit the composed stream before the run, at the cost of one plan per shape rather than one build per shape.
A deployment whose shapes are fixed therefore needs no dispatch-time planner, however large those shapes are, which is what a fixed chunk size secures~\citep{xu2025llmnpu}.
An agent appends whatever its tools return to a context that grows each turn, so its shapes belong to the turn rather than to the deployment, and planning each one in a fraction of a second is what avoids enumerating them in advance.
% SOURCE: the split builds are per configuration and not per shape, 13-14 s from unchanged source
% (Section~\ref{subsec:split}); planning and composing are per shape, 63-206 ms and 56-148 ms on Llama-3.2-1B and
% 0.3-1.1 s and 143-368 ms on Qwen3-4B (Section~\ref{subsec:setup}); the decode design is one build per
% 32,768-token window, with the position setting trip counts at runtime (figs/data/decode-1b.json, windows 32768).
% DECISION: stated as a condition, not a boundary, and the compiler claim is argued rather than demonstrated: no
% pass in our trees composes across operators, so the sentence says "could" and names no gain.


\bibliography{refs}
\bibliographystyle{mlsys2025}

\end{document}
